\documentclass[11pt,a4paper]{article}
\usepackage{fontspec,amsmath,amssymb,amsthm,microtype}
\setmainfont{lmroman10-regular.otf}[BoldFont=lmroman10-bold.otf,ItalicFont=lmroman10-italic.otf,BoldItalicFont=lmroman10-bolditalic.otf]
\usepackage{xcolor}
\usepackage[most]{tcolorbox}
\usepackage{mathrsfs}
\providecommand{\lt}{<}
\providecommand{\gt}{>}
\setmonofont{DejaVuSansMono.ttf}[BoldFont=DejaVuSansMono-Bold.ttf,ItalicFont=DejaVuSansMono-Oblique.ttf,BoldItalicFont=DejaVuSansMono-BoldOblique.ttf,Scale=0.86]
\newfontfamily\leanfallback{STIXTwoMath-Regular.otf}[Scale=0.86]
\definecolor{leanrule}{gray}{0.62}
\definecolor{leanback}{gray}{0.962}
\definecolor{leantitle}{gray}{0.22}
\definecolor{leanmuted}{gray}{0.38}
\newtcolorbox{leanbox}[1]{enhanced,breakable,colback=leanback,colframe=leanrule,
  boxrule=0.4pt,arc=0.6mm,left=2.4mm,right=2.4mm,top=1.8mm,bottom=1.8mm,
  before skip=0.8em,after skip=1em,fonttitle=\small\bfseries,coltitle=leantitle,
  colbacktitle=leanback,titlerule=0.3pt,title={#1}}
\newlength{\leancharwidth}
\newenvironment{leancode}{\par\smallskip\begingroup\ttfamily\footnotesize\raggedright
  \settowidth{\leancharwidth}{M}\setlength{\parskip}{0pt}\setlength{\parindent}{0pt}
  \frenchspacing}{\par\endgroup\smallskip}
\newcommand{\leanline}[2]{\hangindent=\dimexpr#1\leancharwidth+4\leancharwidth\relax
  \hangafter=1\noindent\hspace*{#1\leancharwidth}#2\par}
\newcommand{\leansource}[1]{\par{\footnotesize\color{leanmuted}#1}\par}
\newcommand{\leanlabel}[1]{\par\smallskip\noindent{\small\bfseries #1}\par}
\usepackage[margin=25mm]{geometry}
\usepackage[unicode,hidelinks]{hyperref}
\hypersetup{pdfauthor={LangevinLean Project}}
\setlength{\emergencystretch}{3em}
\setlength{\parskip}{0.4em}
\setlength{\parindent}{0pt}
\allowdisplaybreaks[2]
\raggedbottom
\title{\texorpdfstring{The Gibbs-initialized kinetic Langevin integral solution}{The Gibbs-initialized kinetic Langevin integral solution}}
\author{LangevinLean Project}
\date{}
\begin{document}\maketitle
Draft mathematical exposition; these notes may contain mistakes.


The first-order potential assumptions construct a square-integrable initial phase and an independent specified Brownian source. The selected Picard solution satisfies the kinetic integral equation. The previously proved splitting limit establishes its exact Gaussian-times-Gibbs law at every time; stationarity is a conclusion of this construction.



The initial Gibbs phase and specified Brownian source are realized on an independent product probability space. Splitting convergence and the verified joint source law establish stationarity of the selected kinetic integral solution.


\section{\texorpdfstring{The two target projection moments}{The two target projection moments}}


The target is the product \(\pi=\mathcal N(0,I)\otimes\pi_f\), where \(\pi_f(dx)=Z_f^{-1}e^{-f(x)}dx\). Its velocity projection lies in \(L^{2}\) because a finite-dimensional Gaussian has a finite second moment. Its position projection lies in \(L^{2}\) because strong convexity with \(m>0\) gives Gaussian tail domination for the normalized Gibbs measure. First, use that \(\pi_f\) has total mass one. Factor the phase product and compose the existing Gaussian or Gibbs \(L^{2}\) identity with the corresponding measure-preserving product projection. These lemmas prove the exact integrability required by the random initial state; they do not introduce moment hypotheses.



The target’s product structure gives the position second moment and the exact Gaussian velocity second moment. These ensure that the phase initialization belongs to the square-integrable Hilbert space. They are derived from the target rather than imposed as unexplained moment assumptions.

\begin{leanbox}{Lean: \texttt{kineticGibbsTarget\_\allowbreak{}velocity\_\allowbreak{}memLp}, \texttt{kineticGibbsTarget\_\allowbreak{}position\_\allowbreak{}memLp}}
\leansource{\texttt{kineticGibbsTarget\_\allowbreak{}velocity\_\allowbreak{}memLp} \quad KineticLangevin/\allowbreak{}CanonicalStationarity.lean, lines 21–27}
\begin{leancode}
\leanline{0}{theorem\ kineticGibbsTarget\_velocity\_memLp\ (f\ :\ E\ →\ ℝ)\ (m\ :\ ℝ)\ (hm\ :\ 0<m)}
\leanline{4}{(hf\ :\ Differentiable\ ℝ\ f)\ (hsc\ :\ StrongConvexOn\ univ\ m\ f)\ :}
\leanline{4}{MemLp\ (Prod.fst\ :\ E×E\ →\ E)\ 2\ (kineticGibbsTarget\ f)}
\end{leancode}
\leansource{\texttt{kineticGibbsTarget\_\allowbreak{}position\_\allowbreak{}memLp} \quad KineticLangevin/\allowbreak{}CanonicalStationarity.lean, lines 30–35}
\begin{leancode}
\leanline{0}{theorem\ kineticGibbsTarget\_position\_memLp\ (f\ :\ E\ →\ ℝ)\ (m\ :\ ℝ)\ (hm\ :\ 0<m)}
\leanline{4}{(hf\ :\ Differentiable\ ℝ\ f)\ (hsc\ :\ StrongConvexOn\ univ\ m\ f)\ :}
\leanline{4}{MemLp\ (Prod.snd\ :\ E×E\ →\ E)\ 2\ (kineticGibbsTarget\ f)}
\end{leancode}
\leanlabel{How the Lean proof follows the argument}
\texttt{kinetic\allowbreak{}Gibbs\allowbreak{}Target\_\allowbreak{}velocity\_\allowbreak{}mem\allowbreak{}Lp} gives the formal counterpart of this step.




\texttt{kinetic\allowbreak{}Gibbs\allowbreak{}Target\_\allowbreak{}position\_\allowbreak{}mem\allowbreak{}Lp} gives the formal counterpart of this step.
\end{leanbox}


Earlier results used here: \cite[Section~3.26.13]{LangevinLean2026ULASC}; \cite[Section~3.26.15]{LangevinLean2026ULASC}; \cite[Section~3.26.3]{LangevinLean2026ULASC}.

\section{\texorpdfstring{Build the independent Gibbs and Brownian probability space}{Build the independent Gibbs and Brownian probability space}}


The state space is \(E=\mathbb R^d\), defined as Euclidean space on \(\mathrm{Fin}\,d\). The probability space is the product of the normalized phase law with the already constructed specified Brownian source. Its initial phase is the two \(L^{2}\) coordinate projections, whose integrability was proved above. Thus the starting velocity is standard Gaussian, the starting position has the normalized Gibbs law, and the whole Brownian source is independent of both. The derived projection \(square-integrability\) statements give the two initial \(L^{2}\) coordinates. This fixes the initial state and its independent Brownian forcing before proving invariance.



Take the product of the normalized Gibbs phase measure and the specified Brownian probability space. The coordinate projections have exactly those laws and are independent. This probability space supports both the initial phase and all future filtered responses.

\begin{leanbox}{Lean: \texttt{CanonicalGibbsSpace}, \texttt{canonicalGibbsProbability}, \texttt{canonicalGibbsInitial}}
\leansource{\texttt{CanonicalGibbsSpace} \quad KineticLangevin/\allowbreak{}CanonicalStationarity.lean, lines 39–39}
\begin{leancode}
\leanline{0}{abbrev\ CanonicalGibbsSpace\ (d\ :\ ℕ)\ :=\ EuclideanSpace\ ℝ\ (Fin\ d)}
\end{leancode}
\leansource{\texttt{canonicalGibbsProbability} \quad KineticLangevin/\allowbreak{}CanonicalStationarity.lean, lines 42–43}
\begin{leancode}
\leanline{0}{abbrev\ canonicalGibbsProbability\ (f\ :\ CanonicalGibbsSpace\ d\ →\ ℝ)\ :=}
\leanline{2}{CanonicalProduct.probability\ (kineticGibbsTarget\ f)}
\end{leancode}
\leansource{\texttt{canonicalGibbsInitial} \quad KineticLangevin/\allowbreak{}CanonicalStationarity.lean, lines 46–52}
\begin{leancode}
\leanline{0}{def\ canonicalGibbsInitial\ (f\ :\ CanonicalGibbsSpace\ d\ →\ ℝ)\ (m\ :\ ℝ)\ (hm\ :\ 0<m)}
\leanline{4}{(hf\ :\ Differentiable\ ℝ\ f)\ (hsc\ :\ StrongConvexOn\ univ\ m\ f)\ :}
\leanline{4}{Lp\ (CanonicalGibbsSpace\ d)\ 2\ (canonicalGibbsProbability\ f)\ ×}
\leanline{6}{Lp\ (CanonicalGibbsSpace\ d)\ 2\ (canonicalGibbsProbability\ f)\ :=\ by}
\leanline{2}{haveI\ :=\ kineticGibbsTarget\_probability\ f\ m\ hm\ hf\ hsc}
\leanline{2}{exact\ CanonicalProduct.initial\ (kineticGibbsTarget\ f)}
\leanline{4}{(kineticGibbsTarget\_velocity\_memLp\ f\ m\ hm\ hf\ hsc)\ (kineticGibbsTarget\_position\_memLp\ f\ m\ hm\ hf\ hsc)}
\end{leancode}
\leanlabel{How the Lean proof follows the argument}
\texttt{Canonical\allowbreak{}Gibbs\allowbreak{}Space} names the mathematical object defined above. Its defining expression appears in the adjacent statement; subsequent identities may unfold this definition directly.




\texttt{canonical\allowbreak{}Gibbs\allowbreak{}Probability} names the mathematical object defined above. Its defining expression appears in the adjacent statement; subsequent identities may unfold this definition directly.




\texttt{canonical\allowbreak{}Gibbs\allowbreak{}Initial} names the mathematical object defined above. Its defining expression appears in the adjacent statement; subsequent identities may unfold this definition directly.
\end{leanbox}

\section{\texorpdfstring{The fluctuation-dissipation input}{The fluctuation-dissipation input}}


On the independent product source, the forcing is \(B_t=\sqrt{2\gamma}\,W_t\), where the vector Brownian motion has the established covariance and independent increment laws. It is a continuous function of time taking values in the \(L^{2}\) space, and \(B_0=0\). The proof pulls back the specified Brownian coordinates along the product projection, carries over their proved \(L^{2}\) continuity, and rewrites their value at zero. These are properties of the constructed input, rather than regularity assumptions on an unspecified process. The drift is the pointwise potential gradient lifted to this same \(L^{2}\) space; its global Lipschitz constant is inherited from the stated gradient bound.



The additive velocity input is the vector Brownian process multiplied by \(\sqrt{2\gamma}\). Its filtered increments have thermal Gaussian law required by the splitting. The amplitude is fixed by the fluctuation–dissipation relation, ensuring the target velocity temperature is one.

\begin{leanbox}{Lean: \texttt{canonicalGibbsInput}, \texttt{continuous\_\allowbreak{}canonicalGibbsInput}, \texttt{canonicalGibbsInput\_\allowbreak{}zero} and 1 more}
\leansource{\texttt{canonicalGibbsInput} \quad KineticLangevin/\allowbreak{}CanonicalStationarity.lean, lines 55–59}
\begin{leancode}
\leanline{0}{def\ canonicalGibbsInput\ (f\ :\ CanonicalGibbsSpace\ d\ →\ ℝ)\ (m\ :\ ℝ)\ (hm\ :\ 0<m)}
\leanline{4}{(hf\ :\ Differentiable\ ℝ\ f)\ (hsc\ :\ StrongConvexOn\ univ\ m\ f)\ (γ\ :\ ℝ)\ :}
\leanline{4}{ℝ\ →\ Lp\ (CanonicalGibbsSpace\ d)\ 2\ (canonicalGibbsProbability\ f)\ :=\ by}
\leanline{2}{haveI\ :=\ kineticGibbsTarget\_probability\ f\ m\ hm\ hf\ hsc}
\leanline{2}{exact\ CanonicalProduct.input\ (kineticGibbsTarget\ f)\ γ}
\end{leancode}
\leansource{\texttt{continuous\_\allowbreak{}canonicalGibbsInput} \quad KineticLangevin/\allowbreak{}CanonicalStationarity.lean, lines 61–65}
\begin{leancode}
\leanline{0}{theorem\ continuous\_canonicalGibbsInput\ (f\ :\ CanonicalGibbsSpace\ d\ →\ ℝ)\ (m\ :\ ℝ)\ (hm\ :\ 0<m)}
\leanline{4}{(hf\ :\ Differentiable\ ℝ\ f)\ (hsc\ :\ StrongConvexOn\ univ\ m\ f)\ (γ\ :\ ℝ)\ :}
\leanline{4}{Continuous\ (canonicalGibbsInput\ f\ m\ hm\ hf\ hsc\ γ)}
\end{leancode}
\leansource{\texttt{canonicalGibbsInput\_\allowbreak{}zero} \quad KineticLangevin/\allowbreak{}CanonicalStationarity.lean, lines 67–71}
\begin{leancode}
\leanline{0}{@[simp]\ theorem\ canonicalGibbsInput\_zero\ (f\ :\ CanonicalGibbsSpace\ d\ →\ ℝ)\ (m\ :\ ℝ)\ (hm\ :\ 0<m)}
\leanline{4}{(hf\ :\ Differentiable\ ℝ\ f)\ (hsc\ :\ StrongConvexOn\ univ\ m\ f)\ (γ\ :\ ℝ)\ :}
\leanline{4}{canonicalGibbsInput\ f\ m\ hm\ hf\ hsc\ γ\ 0=0}
\end{leancode}
\leansource{\texttt{canonicalGibbsDrift} \quad KineticLangevin/\allowbreak{}CanonicalStationarity.lean, lines 74–80}
\begin{leancode}
\leanline{0}{def\ canonicalGibbsDrift\ (f\ :\ CanonicalGibbsSpace\ d\ →\ ℝ)\ (m\ :\ ℝ)\ (hm\ :\ 0<m)\ (M\ :\ ℝ≥0)}
\leanline{4}{(hf\ :\ Differentiable\ ℝ\ f)\ (hsc\ :\ StrongConvexOn\ univ\ m\ f)}
\leanline{4}{(hLip\ :\ LipschitzWith\ M\ (gradient\ f))\ :}
\leanline{4}{Lp\ (CanonicalGibbsSpace\ d)\ 2\ (canonicalGibbsProbability\ f)\ →}
\leanline{6}{Lp\ (CanonicalGibbsSpace\ d)\ 2\ (canonicalGibbsProbability\ f)\ :=\ by}
\leanline{2}{haveI\ :=\ kineticGibbsTarget\_probability\ f\ m\ hm\ hf\ hsc}
\leanline{2}{exact\ ULA.lpDrift\ (canonicalGibbsProbability\ f)\ (gradient\ f)\ hLip}
\end{leancode}
\leanlabel{How the Lean proof follows the argument}
\texttt{canonical\allowbreak{}Gibbs\allowbreak{}Input} names the mathematical object defined above. Its defining expression appears in the adjacent statement; subsequent identities may unfold this definition directly.




\texttt{continuous\_\allowbreak{}canonical\allowbreak{}Gibbs\allowbreak{}Input} gives the formal counterpart of this step. Named proof applications include \texttt{Canonical\allowbreak{}Product.\allowbreak{}continuous\_\allowbreak{}input}; the adjacent proof shows their hypotheses and order of use.




\texttt{canonical\allowbreak{}Gibbs\allowbreak{}Input\_\allowbreak{}zero} gives the formal counterpart of this step. Named proof applications include \texttt{Canonical\allowbreak{}Product.\allowbreak{}input\_\allowbreak{}zero}; the adjacent proof shows their hypotheses and order of use.




\texttt{canonical\allowbreak{}Gibbs\allowbreak{}Drift} names the mathematical object defined above. Its defining expression appears in the adjacent statement; subsequent identities may unfold this definition directly.
\end{leanbox}


Earlier results used here: \cite[Section~3.13.2]{LangevinLean2026ULASC}.

\section{\texorpdfstring{Select a continuous solution and record its equation}{Select a continuous solution and record its equation}}


For a horizon \(T\ge0\), Picard existence for a globally Lipschitz additive equation constructs the phase path \(Y_t=(V_t,X_t)\) in \(L^{2}\). It satisfies \[V_t-V_0=-\int_0^t(\gamma V_s+\nabla f(X_s))\,ds+\sqrt{2\gamma}\,W_t,\qquad X_t-X_0=\int_0^t V_s\,ds.\] Product notation combines these two equations into one integral identity using the negative drift. The flow definition supplies the proved \(L^{2}\) Lipschitz gradient, the constructed initial state, and the continuous input to the restarted kinetic solution. Its specification theorem then proves continuity, the exact initial value, and the displayed integral equation on the whole interval. No stationary process or solution identity is postulated.



Existence and uniqueness give a continuous square-integrable phase path with the prescribed initial state and input. Recording its integral equation links this object to the physical kinetic dynamics. Stationarity is then a property to prove about that specific path.

\begin{leanbox}{Lean: \texttt{canonicalGibbsFlow}, \texttt{canonicalGibbsFlow\_\allowbreak{}spec}}
\leansource{\texttt{canonicalGibbsFlow} \quad KineticLangevin/\allowbreak{}CanonicalStationarity.lean, lines 83–92}
\begin{leancode}
\leanline{0}{def\ canonicalGibbsFlow\ (f\ :\ CanonicalGibbsSpace\ d\ →\ ℝ)\ (m\ :\ ℝ)\ (hm\ :\ 0<m)\ (M\ γ\ :\ ℝ≥0)}
\leanline{4}{(hf\ :\ Differentiable\ ℝ\ f)\ (hsc\ :\ StrongConvexOn\ univ\ m\ f)}
\leanline{4}{(hLip\ :\ LipschitzWith\ M\ (gradient\ f))\ (T\ :\ ℝ)\ (hT\ :\ 0≤T)\ :}
\leanline{4}{ℝ\ →\ Lp\ (CanonicalGibbsSpace\ d)\ 2\ (canonicalGibbsProbability\ f)\ ×}
\leanline{6}{Lp\ (CanonicalGibbsSpace\ d)\ 2\ (canonicalGibbsProbability\ f)\ :=\ by}
\leanline{2}{haveI\ :=\ kineticGibbsTarget\_probability\ f\ m\ hm\ hf\ hsc}
\leanline{2}{exact\ kineticSegmentFlow\ (ULA.lpDrift\ (canonicalGibbsProbability\ f)\ (gradient\ f)\ hLip)\ M\ γ\ 1}
\leanline{4}{(ULA.lpDrift\_lipschitz\ (canonicalGibbsProbability\ f)\ (gradient\ f)\ hLip)\ T\ hT}
\leanline{4}{(canonicalGibbsInitial\ f\ m\ hm\ hf\ hsc)\ (canonicalGibbsInput\ f\ m\ hm\ hf\ hsc\ γ)}
\leanline{4}{(continuous\_canonicalGibbsInput\ f\ m\ hm\ hf\ hsc\ γ)}
\end{leancode}
\leansource{\texttt{canonicalGibbsFlow\_\allowbreak{}spec} \quad KineticLangevin/\allowbreak{}CanonicalStationarity.lean, lines 95–110}
\begin{leancode}
\leanline{0}{theorem\ canonicalGibbsFlow\_spec\ (f\ :\ CanonicalGibbsSpace\ d\ →\ ℝ)\ (m\ :\ ℝ)\ (hm\ :\ 0<m)\ (M\ γ\ :\ ℝ≥0)}
\leanline{4}{(hf\ :\ Differentiable\ ℝ\ f)\ (hsc\ :\ StrongConvexOn\ univ\ m\ f)}
\leanline{4}{(hLip\ :\ LipschitzWith\ M\ (gradient\ f))\ (T\ :\ ℝ)\ (hT\ :\ 0≤T)\ :}
\leanline{4}{Continuous\ (canonicalGibbsFlow\ f\ m\ hm\ M\ γ\ hf\ hsc\ hLip\ T\ hT)\ ∧}
\leanline{4}{canonicalGibbsFlow\ f\ m\ hm\ M\ γ\ hf\ hsc\ hLip\ T\ hT\ 0=canonicalGibbsInitial\ f\ m\ hm\ hf\ hsc\ ∧}
\leanline{4}{∀\ t\ ∈\ Icc\ 0\ T,}
\leanline{6}{canonicalGibbsFlow\ f\ m\ hm\ M\ γ\ hf\ hsc\ hLip\ T\ hT\ t-canonicalGibbsInitial\ f\ m\ hm\ hf\ hsc\ =}
\leanline{8}{-(∫\ s\ in\ 0..t,\ kineticNegativeDrift}
\leanline{10}{(canonicalGibbsDrift\ f\ m\ hm\ M\ hf\ hsc\ hLip)\ γ\ 1}
\leanline{10}{(canonicalGibbsFlow\ f\ m\ hm\ M\ γ\ hf\ hsc\ hLip\ T\ hT\ s))+}
\leanline{8}{(canonicalGibbsInput\ f\ m\ hm\ hf\ hsc\ γ\ t,0)}
\end{leancode}
\leanlabel{How the Lean proof follows the argument}
\texttt{canonical\allowbreak{}Gibbs\allowbreak{}Flow} names the mathematical object defined above. Its defining expression appears in the adjacent statement; subsequent identities may unfold this definition directly.




\texttt{canonical\allowbreak{}Gibbs\allowbreak{}Flow\_\allowbreak{}spec} gives the formal counterpart of this step. Named proof applications include \texttt{kinetic\allowbreak{}Segment\allowbreak{}Flow\_\allowbreak{}spec}; the adjacent proof shows their hypotheses and order of use.
\end{leanbox}


Earlier results used here: \cite[Section~3.13.2]{LangevinLean2026ULASC}; \cite[Section~3.13.5]{LangevinLean2026ULASC}.

\section{\texorpdfstring{The exact target law at every time}{The exact target law at every time}}


Fix \(t\in[0,T]\). On every mesh size \(h\ge0\), the initial phase and the entire filtered Brownian innovation stream have joint law \[\pi\otimes Q_h^{\otimes\mathbb N},\qquad Q_h=\mathcal N\!\left(0,(1-e^{-2\gamma h})I\right).\] The argument obtains this from the proved product-source joint-law theorem and the exact identity between filtered Brownian increments and OU innovations. It restricts the solution equation from \([0,T]\) to \([0,t]\), then applies the splitting stationarity theorem. That theorem has already proved convergence of the splitting chain, an order-\(h^2\) stationary weak defect from Metropolis rejection, Markov telescoping, and continuity of the limiting law. The conclusion is the exact pushforward identity \[\operatorname{Law}(V_t,X_t)=\mathcal N(0,I)\otimes\pi_f.\] Every source, initial moment, solution, innovation law, and limit used here has been constructed or proved; there is no invariant-target or output-law assumption remaining in this conclusion.



Apply the splitting-limit stationarity transfer at an arbitrary time in the horizon. The verified joint input law and the exact splitting convergence verify its hypotheses. The resulting phase law is the full product target, yielding both stationary marginals simultaneously.

\begin{leanbox}{Lean: \texttt{canonicalGibbsFlow\_\allowbreak{}stationary}}
\leansource{\texttt{canonicalGibbsFlow\_\allowbreak{}stationary} \quad KineticLangevin/\allowbreak{}CanonicalStationarity.lean, lines 114–136}
\begin{leancode}
\leanline{0}{theorem\ canonicalGibbsFlow\_stationary\ (f\ :\ CanonicalGibbsSpace\ d\ →\ ℝ)\ (m\ :\ ℝ)\ (hm\ :\ 0<m)\ (M\ γ\ :\ ℝ≥0)}
\leanline{4}{(hf\ :\ Differentiable\ ℝ\ f)\ (hsc\ :\ StrongConvexOn\ univ\ m\ f)}
\leanline{4}{(hLip\ :\ LipschitzWith\ M\ (gradient\ f))\ (T\ :\ ℝ)\ (hT\ :\ 0≤T)\ (t\ :\ ℝ)\ (ht\ :\ t∈Icc\ 0\ T)\ :}
\leanline{4}{(canonicalGibbsProbability\ f).map\ (lpPhaseCoe\ (canonicalGibbsProbability\ f)}
\leanline{6}{(canonicalGibbsFlow\ f\ m\ hm\ M\ γ\ hf\ hsc\ hLip\ T\ hT\ t))=kineticGibbsTarget\ f}
\end{leancode}
\leanlabel{How the Lean proof follows the argument}
\texttt{canonical\allowbreak{}Gibbs\allowbreak{}Flow\_\allowbreak{}stationary} gives the formal counterpart of this step. The main intermediate facts are \texttt{hs}, \texttt{hjoint} in that order. Named proof applications include \texttt{kinetic\_\allowbreak{}integral\_\allowbreak{}stationary\_\allowbreak{}of\_\allowbreak{}joint\allowbreak{}Noise}, \texttt{hs.\allowbreak{}2.\allowbreak{}2}; the adjacent proof shows their hypotheses and order of use.
\end{leanbox}

\begin{thebibliography}{ULA-LSI}
\bibitem[ULA-SC]{LangevinLean2026ULASC}
LangevinLean Project. \emph{Unadjusted Langevin under Strong Convexity}. LL-01, 2026. Draft report, content revision d34ead06987c. \url{https://odakyildiz.github.io/LangevinLean/docs/pdf/UnadjustedLangevin/FullProof.pdf}.
\end{thebibliography}

\end{document}
